% Згенеровано tools/buildbib.py -- не редагувати вручну

\bibitem{ch15:notpetya-ua} Аналіз атаки NotPetya як прикладу компрометації ланцюга постачання програмного забезпечення в Україні. \emph{Захист інформації}. 2023. Т.~25, №~1.

\bibitem{ch12:apt-ua} Аналіз діяльності APT-угруповань, спрямованих проти критичної інфраструктури України. \emph{Захист інформації}. 2023. Т.~25, №~2.

\bibitem{ch01:apt-ua} Аналіз кібератак та діяльності APT груп в Україні. \emph{Кібербезпека: освіта, наука, техніка}. 2024. URL: \url{https://csecurity.kubg.edu.ua/index.php/journal/article/view/574}.

\bibitem{ch10:edr-ua} Архітектура та можливості систем виявлення й реагування на кінцевих точках (EDR) у корпоративних мережах. \emph{Кібербезпека: освіта, наука, техніка}. 2024. URL: \url{https://csecurity.kubg.edu.ua/index.php/journal/article/view/731}.

\bibitem{ch08:packer-ua} Виявлення запакованого шкідливого програмного забезпечення на основі аналізу ентропії та структурних ознак. \emph{Захист інформації}. 2023. Т.~25, №~3.

\bibitem{ch13:apt-crit-ua} Виявлення цілеспрямованих атак на об'єкти критичної інфраструктури України: методи та підходи. \emph{Захист інформації}. 2023. Т.~25, №~3.

\bibitem{ch05:windows-vuln} Дослідження методів виявлення вразливостей програмного забезпечення та ІКС на базі ОС Windows із розробленням засобу автоматизованого аналізу. \emph{Кібербезпека: освіта, наука, техніка}. 2026. URL: \url{https://csecurity.kubg.edu.ua/index.php/journal/article/view/1302}.

\bibitem{dstu-27001} ДСТУ EN ISO/IEC 27001:2023. Інформаційна безпека, кібербезпека та захист приватності. Системи керування інформаційною безпекою. Вимоги. Київ: ДП «УкрНДНЦ», 2023.

\bibitem{dstu-27002} ДСТУ EN ISO/IEC 27002:2023. Інформаційна безпека, кібербезпека та захист приватності. Заходи інформаційної безпеки. Київ: ДП «УкрНДНЦ», 2023.

\bibitem{kmu-518} Загальні вимоги до кіберзахисту об'єктів критичної інфраструктури: затв. постановою Кабінету Міністрів України від 19.06.2019 № 518. URL: \url{https://zakon.rada.gov.ua/laws/show/518-2019-п}.

\bibitem{ch15:supplychain-ua} Загрози ланцюга постачання програмного забезпечення та методи забезпечення цілісності й довіри до компонентів. \emph{Кібербезпека: освіта, наука, техніка}. 2024. URL: \url{https://csecurity.kubg.edu.ua/index.php/journal/article/view/802}.

\bibitem{ch07:dynamic-ua} Застосування динамічного аналізу та пісочниць для дослідження шкідливого програмного забезпечення. \emph{Кібербезпека: освіта, наука, техніка}. 2024. URL: \url{https://csecurity.kubg.edu.ua/index.php/journal/article/view/615}.

\bibitem{ch09:ml-ua} Застосування методів машинного навчання для виявлення шкідливого програмного забезпечення. \emph{Кібербезпека: освіта, наука, техніка}. 2024. URL: \url{https://csecurity.kubg.edu.ua/index.php/journal/article/view/688}.

\bibitem{ch12:ti-ua} Застосування розвідки кіберзагроз (threat intelligence) для проактивного захисту інформаційних систем. \emph{Кібербезпека: освіта, наука, техніка}. 2024. URL: \url{https://csecurity.kubg.edu.ua/index.php/journal/article/view/759}.

\bibitem{ch14:ics-ua} Захист автоматизованих систем керування технологічними процесами об'єктів критичної інфраструктури України від кіберзагроз. \emph{Кібербезпека: освіта, наука, техніка}. 2024. URL: \url{https://csecurity.kubg.edu.ua/index.php/journal/article/view/788}.

\bibitem{ch04:ilienko} Ільєнко А. В., Ільєнко С. С., Яковенко О. Л., Галич Є. О., Павленко В. Г. Застосування штучного інтелекту для прогнозування кіберзагроз. \emph{Кібербезпека: освіта, наука, техніка}. 2024. URL: \url{https://csecurity.kubg.edu.ua/index.php/journal/article/view/631}.

\bibitem{cert-ua} Команда реагування на комп'ютерні надзвичайні події України CERT-UA. Аналітичні звіти та індикатори компрометації. Державна служба спеціального зв'язку та захисту інформації України. URL: \url{https://cert.gov.ua/}.

\bibitem{kk-ukr} Кримінальний кодекс України: статті 361 (несанкціоноване втручання в роботу ЕОМ), 361-1 (створення з метою використання, поширення або збуту шкідливих програмних чи технічних засобів), 361-2, 362, 363, 363-1. URL: \url{https://zakon.rada.gov.ua/laws/show/2341-14}.

\bibitem{ch02:kpk} Кримінальний процесуальний кодекс України від 13.04.2012 № 4651-VI (щодо вилучення й дослідження електронних доказів). URL: \url{https://zakon.rada.gov.ua/laws/show/4651-17}.

\bibitem{ch04:lungol} Лунгол О. М. Багаторівнева модель атрибуції кібератак. \emph{Кібербезпека: освіта, наука, техніка}. 2025. Спеціальний випуск. URL: \url{https://csecurity.kubg.edu.ua/index.php/journal/article/view/904}.

\bibitem{ch08:obfusc-ua} Методи обфускації програмного коду та підходи до їх подолання в задачах аналізу шкідливого програмного забезпечення. \emph{Кібербезпека: освіта, наука, техніка}. 2024. URL: \url{https://csecurity.kubg.edu.ua/index.php/journal/article/view/702}.

\bibitem{ch16:forensic-ua} Методи цифрової криміналістики при розслідуванні інцидентів, пов'язаних зі шкідливим програмним забезпеченням. \emph{Захист інформації}. 2023. Т.~25, №~2.

\bibitem{ch03:hybrid} Метод захисту програмного забезпечення на основі гібридного аналізу коду. \emph{Кібербезпека: освіта, наука, техніка}. 2025. URL: \url{https://csecurity.kubg.edu.ua/index.php/journal/article/view/871}.

\bibitem{ch14:iot-ua} Моделі та методи виявлення шкідливої активності у середовищах Інтернету речей. \emph{Захист інформації}. 2023. Т.~25, №~4.

\bibitem{ch10:monitor-ua} Моніторинг подій безпеки та виявлення аномальної активності на кінцевих вузлах. \emph{Захист інформації}. 2023. Т.~25, №~4.

\bibitem{ndtzi-2-5-004} НД ТЗІ 2.5-004-99. Критерії оцінки захищеності інформації в комп'ютерних системах від несанкціонованого доступу. Київ: ДСТСЗІ СБ України, 1999.

\bibitem{ndtzi-3-7-003} НД ТЗІ 3.7-003-05. Порядок проведення робіт із створення комплексної системи захисту інформації в інформаційно-телекомунікаційній системі. Київ: ДСТСЗІ СБ України, 2005.

\bibitem{ndtzi-1-1-003} НД ТЗІ 1.1-003-99. Термінологія в галузі захисту інформації в комп'ютерних системах від несанкціонованого доступу. Київ: ДСТСЗІ СБ України, 1999. URL: \url{https://cip.gov.ua/ua/news/normativni-dokumenti-z-pitan-tzi}.

\bibitem{ch02:ndtzi-1-4} НД ТЗІ 1.4-001-2000. Типове положення про службу захисту інформації в автоматизованій системі. Київ: ДСТСЗІ СБ України, 2000.

\bibitem{ch09:neural-ua} Нейромережеві моделі класифікації шкідливого програмного забезпечення за статичними ознаками. \emph{Захист інформації}. 2024. Т.~26, №~1.

\bibitem{ch03:opirskyi-ml} Опірський І. Р., Мельничук М. Р. Застосування машинного навчання в статичному аналізі безпеки коду. \emph{Кібербезпека: освіта, наука, техніка}. 2026. URL: \url{https://csecurity.kubg.edu.ua/index.php/journal/article/view/1193}.

\bibitem{ch17:governance-ua} Організаційно-управлінські засади побудови системи кіберзахисту підприємства в умовах сучасних загроз. \emph{Кібербезпека: освіта, наука, техніка}. 2024. URL: \url{https://csecurity.kubg.edu.ua/index.php/journal/article/view/830}.

\bibitem{ch16:ir-ua} Організація процесу реагування на кіберінциденти в установах України відповідно до вимог законодавства. \emph{Кібербезпека: освіта, наука, техніка}. 2024. URL: \url{https://csecurity.kubg.edu.ua/index.php/journal/article/view/816}.

\bibitem{ch07:behavior-ua} Поведінкові методи виявлення шкідливого програмного забезпечення на основі аналізу системних викликів. \emph{Захист інформації}. 2023. Т.~25, №~2.

\bibitem{ch02:polotai-ids} Полотай О. І. та ін. Інтеграція IDS/IPS у корпоративні мережі та перевірка playbook на кіберполігонах. \emph{Кібербезпека: освіта, наука, техніка}. 2025. URL: \url{https://csecurity.kubg.edu.ua/index.php/journal/article/view/889}.

\bibitem{ch13:hunt-ua} Проактивне полювання за кіберзагрозами (threat hunting) як складова сучасної стратегії кіберзахисту. \emph{Кібербезпека: освіта, наука, техніка}. 2024. URL: \url{https://csecurity.kubg.edu.ua/index.php/journal/article/view/773}.

\bibitem{ch02:zu-ek} Про електронні комунікації: Закон України від 16.12.2020 № 1089-IX. URL: \url{https://zakon.rada.gov.ua/laws/show/1089-20}.

\bibitem{cip} Про затвердження Загальних вимог до кіберзахисту об'єктів критичної інфраструктури та Порядку їх застосування: постанова Кабінету Міністрів України. URL: \url{https://cip.gov.ua/}.

\bibitem{zu-zi} Про захист інформації в інформаційно-комунікаційних системах: Закон України від 05.07.1994 № 80/94-ВР. URL: \url{https://zakon.rada.gov.ua/laws/show/80/94-вр}.

\bibitem{zu-ki} Про критичну інфраструктуру: Закон України від 16.11.2021 № 1882-IX. URL: \url{https://zakon.rada.gov.ua/laws/show/1882-20}.

\bibitem{zu-cyber} Про основні засади забезпечення кібербезпеки України: Закон України від 05.10.2017 № 2163-VIII (редакція від 19.10.2025). URL: \url{https://zakon.rada.gov.ua/laws/show/2163-19}.

\bibitem{ch17:strategy-ua} Про рішення Ради національної безпеки і оборони України «Про Стратегію кібербезпеки України». Указ Президента України. Київ, 2021.

\bibitem{ukaz-447} Про Стратегію кібербезпеки України: Указ Президента України від 26.08.2021 № 447/2021. URL: \url{https://zakon.rada.gov.ua/laws/show/447/2021}.

\bibitem{ch11:detection-ua} Розроблення правил виявлення кіберзагроз для систем керування подіями безпеки (SIEM). \emph{Кібербезпека: освіта, наука, техніка}. 2024. URL: \url{https://csecurity.kubg.edu.ua/index.php/journal/article/view/745}.

\bibitem{ch06:sast} Трофименко О. Г., Дика А. І., Лобода Ю. Г. Порівняльний аналіз методів тестування безпеки застосунків (SAST, DAST, IAST). \emph{Кібербезпека: освіта, наука, техніка}. 2023. Т.~4, №~20. URL: \url{https://csecurity.kubg.edu.ua/index.php/journal/article/view/467}.

\bibitem{ch02:assoc} Угода про асоціацію між Україною та Європейським Союзом: ратифіковано Законом України від 16.09.2014 № 1678-VII. URL: \url{https://zakon.rada.gov.ua/laws/show/984_011}.

\bibitem{ch11:ioc-ua} Формування та застосування індикаторів компрометації у процесах виявлення кіберінцидентів. \emph{Захист інформації}. 2023. Т.~25, №~1.

\bibitem{ch01:shevchenko} Шевченко С. М., Жданова Ю. Д., Гаркушенко А. М. Аналіз сучасних загроз інформаційній безпеці та підходів до їх класифікації. \emph{Кібербезпека: освіта, наука, техніка}. 2025. URL: \url{https://csecurity.kubg.edu.ua/index.php/journal/article/view/945}.

\bibitem{ch09:anderson-gym} Anderson H. S., Kharkar A., Filar B., Evans D., Roth P. Learning to Evade Static PE Machine Learning Malware Models via Reinforcement Learning. \emph{arXiv preprint arXiv:1801.08917}. 2018.

\bibitem{ch09:ember} Anderson H. S., Roth P. EMBER: An Open Dataset for Training Static PE Malware Machine Learning Models. \emph{arXiv preprint arXiv:1804.04637}. 2018.

\bibitem{anderson} Anderson R. Security Engineering: A Guide to Building Dependable Distributed Systems. 3rd ed. Indianapolis: Wiley, 2020. 1232~p. URL: \url{https://www.cl.cam.ac.uk/~rja14/book.html}.

\bibitem{ch17:anderson} Anderson R., Moore T. The Economics of Information Security. \emph{Science}. 2006. Vol.~314, №~5799. P.~610--613. DOI: 10.1126/science.1130992.

\bibitem{ch06:andriesse} Andriesse D., Chen X., van der Veen V., Slowinska A., Bos H. An In-Depth Analysis of Disassembly on Full-Scale x86/x64 Binaries. \emph{Proceedings of the 25th USENIX Security Symposium}. USENIX Association, 2016. P.~583--600.

\bibitem{ch14:antonakakis} Antonakakis M., April T., Bailey M., et al. Understanding the Mirai Botnet. \emph{Proceedings of the 26th USENIX Security Symposium}. USENIX Association, 2017. P.~1093--1110.

\bibitem{ch01:avtest} AV-TEST Institute. Malware Statistics \& Trends Report. Magdeburg: AV-TEST GmbH, 2024. URL: \url{https://www.av-test.org/en/statistics/malware/}.

\bibitem{ch15:ccleaner} Avast Threat Intelligence Team. Progress on CCleaner Investigation. Avast Blog, 2017.

\bibitem{ch03:axelsson} Axelsson S. The Base-Rate Fallacy and the Difficulty of Intrusion Detection. \emph{ACM Transactions on Information and System Security}. 2000. Vol.~3, No.~3. P.~186--205. DOI: 10.1145/357830.357849.

\bibitem{aycock} Aycock J. Computer Viruses and Malware. New York: Springer, 2006. 227~p. (Advances in Information Security, Vol.~22).

\bibitem{ch17:awareness} Bada M., Sasse A. M., Nurse J. R. C. Cyber Security Awareness Campaigns: Why Do They Fail to Change Behaviour? \emph{International Conference on Cyber Security for Sustainable Society}. 2015. P.~118--131.

\bibitem{ch07:bailey} Bailey M., Oberheide J., Andersen J., Mao Z. M., Jahanian F., Nazario J. Automated Classification and Analysis of Internet Malware. \emph{Recent Advances in Intrusion Detection (RAID 2007)}. LNCS Vol.~4637. Berlin: Springer, 2007. P.~178--197.

\bibitem{ch12:attribution} Bartholomew B., Guerrero-Saade J. A. Wave Your False Flags! Deception Tactics Muddying Attribution in Targeted Attacks. \emph{Virus Bulletin Conference}. Denver, 2016.

\bibitem{ch07:bayer} Bayer U., Comparetti P. M., Hlauschek C., Kruegel C., Kirda E. Scalable, Behavior-Based Malware Clustering. \emph{Proceedings of the 16th Annual Network and Distributed System Security Symposium (NDSS)}. San Diego: ISOC, 2009.

\bibitem{pyramid} Bianco D. The Pyramid of Pain. Enterprise Detection \& Response Blog, 2013. URL: \url{http://detect-respond.blogspot.com/2013/03/the-pyramid-of-pain.html}.

\bibitem{ch09:biggio} Biggio B., Roli F. Wild Patterns: Ten Years After the Rise of Adversarial Machine Learning. \emph{Pattern Recognition}. 2018. Vol.~84. P.~317--331.

\bibitem{ch05:blunden} Blunden B. The Rootkit Arsenal: Escape and Evasion in the Dark Corners of the System. 2nd ed. Burlington: Jones \& Bartlett Learning, 2012. 784~p.

\bibitem{ch11:powershell} Bohannon D., Holmes L. Revoke-Obfuscation: PowerShell Obfuscation Detection Using Science. \emph{Black Hat USA}. Las Vegas, 2017.

\bibitem{ch08:branco} Branco R. R., Barbosa G. N., Neto P. D. Scientific but Not Academical Overview of Malware Anti-Debugging, Anti-Disassembly and Anti-VM Technologies. \emph{Black Hat USA}. Las Vegas, 2012.

\bibitem{ch16:orderofvolatility} Brezinski D., Killalea T. Guidelines for Evidence Collection and Archiving. RFC~3227. IETF, 2002.

\bibitem{ch06:brumley} Brumley D., Jager I., Avgerinos T., Schwartz E. J. BAP: A Binary Analysis Platform. \emph{Computer Aided Verification (CAV 2011)}. LNCS Vol.~6806. Berlin: Springer, 2011. P.~463--469.

\bibitem{ch12:diamond} Caltagirone S., Pendergast A., Betz C. The Diamond Model of Intrusion Analysis. Center for Cyber Intelligence Analysis and Threat Research, 2013.

\bibitem{capec} CAPEC: Common Attack Pattern Enumeration and Classification. The MITRE Corporation. URL: \url{https://capec.mitre.org/}.

\bibitem{ch16:carrier} Carrier B. File System Forensic Analysis. Boston: Addison-Wesley, 2005. 600~p.

\bibitem{ch16:casey} Casey E. Digital Evidence and Computer Crime: Forensic Science, Computers, and the Internet. 3rd ed. Waltham: Academic Press, 2011. 840~p.

\bibitem{ch01:cert-report} CERT-UA. Огляд діяльності урядової команди реагування на комп'ютерні надзвичайні події України: щорічні звіти. Київ: Держспецзв'язку, 2022--2024. URL: \url{https://cert.gov.ua/reports}.

\bibitem{ch07:chen} Chen X., Andersen J., Mao Z. M., Bailey M., Nazario J. Towards an Understanding of Anti-Virtualization and Anti-Debugging Behavior in Modern Malware. \emph{IEEE International Conference on Dependable Systems and Networks (DSN)}. Anchorage: IEEE, 2008. P.~177--186.

\bibitem{ch14:industroyer} Cherepanov A. Win32/Industroyer: A New Threat for Industrial Control Systems. ESET White Paper, 2017.

\bibitem{industroyer} Cherepanov A., Lipovsky R. Industroyer: Biggest Threat to Industrial Control Systems since Stuxnet. ESET White Paper. Bratislava: ESET, 2017.

\bibitem{ch03:chess} Chess D. M., White S. R. An Undetectable Computer Virus. \emph{Proceedings of the Virus Bulletin Conference}. 2000.

\bibitem{ch06:christodorescu} Christodorescu M., Jha S., Seshia S. A., Song D., Bryant R. E. Semantics-Aware Malware Detection. \emph{2005 IEEE Symposium on Security and Privacy}. Oakland: IEEE, 2005. P.~32--46. DOI: 10.1109/SP.2005.20.

\bibitem{ch10:gartner-edr} Chuvakin A. Named: Endpoint Threat Detection \& Response. Gartner Blog Network, 2013. URL: \url{https://blogs.gartner.com/anton-chuvakin/2013/07/26/named-endpoint-threat-detection-response/}.

\bibitem{nist-61} Cichonski P., Millar T., Grance T., Scarfone K. Computer Security Incident Handling Guide. NIST Special Publication 800-61, Revision~2. Gaithersburg: NIST, 2012. DOI: 10.6028/NIST.SP.800-61r2.

\bibitem{ch02:cis} CIS Critical Security Controls, Version 8.1. Center for Internet Security, 2024. URL: \url{https://www.cisecurity.org/controls/v8-1}.

\bibitem{ch15:solarwinds} CISA. Advanced Persistent Threat Compromise of Government Agencies, Critical Infrastructure, and Private Sector Organizations (SolarWinds). Cybersecurity and Infrastructure Security Agency Alert AA20-352A, 2020.

\bibitem{ch06:clamav} Cisco-Talos. ClamAV: Open Source Antivirus Engine --- Signatures and Database Formats. URL: \url{https://docs.clamav.net/}.

\bibitem{ch01:cohen-thesis} Cohen F. Computer Viruses. PhD Dissertation. University of Southern California, 1986.

\bibitem{ch01:cohen} Cohen F. Computer Viruses: Theory and Experiments. \emph{Computers \& Security}. 1987. Vol.~6, No.~1. P.~22--35. DOI: 10.1016/0167-4048(87)90122-2.

\bibitem{ch08:collberg} Collberg C., Nagra J. Surreptitious Software: Obfuscation, Watermarking, and Tamperproofing for Software Protection. Boston: Addison-Wesley, 2009. 792~p.

\bibitem{budapest} Convention on Cybercrime (Budapest Convention). ETS No.~185. Budapest: Council of Europe, 2001. URL: \url{https://www.coe.int/en/web/cybercrime/the-budapest-convention}.

\bibitem{ch07:cuckoo} Cuckoo Foundation. Cuckoo Sandbox: Automated Malware Analysis --- Documentation. URL: \url{https://cuckoo.readthedocs.io/}.

\bibitem{cve} CVE: Common Vulnerabilities and Exposures. The MITRE Corporation. URL: \url{https://www.cve.org/}.

\bibitem{nis2} Directive (EU) 2022/2555 of the European Parliament and of the Council of 14 December 2022 on measures for a high common level of cybersecurity across the Union (NIS 2 Directive). \emph{Official Journal of the EU}. L 333. P.~80--152. URL: \url{https://eur-lex.europa.eu/eli/dir/2022/2555/oj}.

\bibitem{wipers2022} Dobrovolsky G. et al. (CERT-UA, ESET, Microsoft). WhisperGate, HermeticWiper, IsaacWiper, CaddyWiper and the 2022 destructive campaigns against Ukraine: joint analyses. 2022.

\bibitem{ch06:bindiff} Dullien T., Rolles R. Graph-based comparison of Executable Objects. \emph{Symposium sur la sécurité des technologies de l'information et des communications (SSTIC)}. 2005.

\bibitem{ch05:eagle} Eagle C. The IDA Pro Book. 2nd ed. San Francisco: No Starch Press, 2011. 672~p.

\bibitem{ch07:egele} Egele M., Scholte T., Kirda E., Kruegel C. A Survey on Automated Dynamic Malware-Analysis Techniques and Tools. \emph{ACM Computing Surveys}. 2012. Vol.~44, №~2. Art.~6. P.~1--42. DOI: 10.1145/2089125.2089126.

\bibitem{ch14:enck} Enck W., Octeau D., McDaniel P., Chaudhuri S. A Study of Android Application Security. \emph{Proceedings of the 20th USENIX Security Symposium}. USENIX Association, 2011. P.~21--36.

\bibitem{enisa} ENISA Threat Landscape 2024. Athens: European Union Agency for Cybersecurity, 2024. URL: \url{https://www.enisa.europa.eu/publications/enisa-threat-landscape-2024}.

\bibitem{ch15:enisa} ENISA. Threat Landscape for Supply Chain Attacks. European Union Agency for Cybersecurity, 2021.

\bibitem{notpetya} ESET Research. New TeleBots backdoor: First evidence linking Industroyer to NotPetya. ESET WeLiveSecurity, 2018. URL: \url{https://www.welivesecurity.com/}.

\bibitem{stuxnet} Falliere N., Murchu L. O., Chien E. W32.Stuxnet Dossier. Version~1.4. Symantec Security Response, 2011.

\bibitem{ch03:fawcett} Fawcett T. An Introduction to ROC Analysis. \emph{Pattern Recognition Letters}. 2006. Vol.~27, No.~8. P.~861--874. DOI: 10.1016/j.patrec.2005.10.010.

\bibitem{ch02:cvss} FIRST. Common Vulnerability Scoring System Version 4.0: Specification Document. Forum of Incident Response and Security Teams, 2023. URL: \url{https://www.first.org/cvss/v4.0/}.

\bibitem{ch04:epss} FIRST. Exploit Prediction Scoring System (EPSS). Forum of Incident Response and Security Teams. URL: \url{https://www.first.org/epss/}.

\bibitem{ch04:fair} Freund J., Jones J. Measuring and Managing Information Risk: A FAIR Approach. Oxford: Butterworth-Heinemann, 2014. 408~p.

\bibitem{ch17:fair} Freund J., Jones J. Measuring and Managing Information Risk: A FAIR Approach. Waltham: Butterworth-Heinemann, 2014. 408~p.

\bibitem{ch16:antiforensics} Garfinkel S. Anti-Forensics: Techniques, Detection and Countermeasures. \emph{Proceedings of the 2nd International Conference on i-Warfare and Security}. 2007. P.~77--84.

\bibitem{ch10:xdr} Gartner. Market Guide for Extended Detection and Response. Gartner Research, 2021.

\bibitem{ch09:gibert} Gibert D., Mateu C., Planes J. The Rise of Machine Learning for Detection and Classification of Malware: Research Developments, Trends and Challenges. \emph{Journal of Network and Computer Applications}. 2020. Vol.~153. Art.~102526.

\bibitem{ch14:android} Google. Android Security Model and Application Sandbox --- Android Open Source Project Documentation. 2024. URL: \url{https://source.android.com/docs/security}.

\bibitem{ch04:gordonloeb} Gordon L. A., Loeb M. P. The Economics of Information Security Investment. \emph{ACM Transactions on Information and System Security}. 2002. Vol.~5, No.~4. P.~438--457. DOI: 10.1145/581271.581274.

\bibitem{ch17:gordon-loeb} Gordon L. A., Loeb M. P. The Economics of Information Security Investment. \emph{ACM Transactions on Information and System Security}. 2002. Vol.~5, №~4. P.~438--457.

\bibitem{ch01:greenberg} Greenberg A. Sandworm: A New Era of Cyberwar and the Hunt for the Kremlin's Most Dangerous Hackers. New York: Doubleday, 2019. 348~p.

\bibitem{ch09:grosse} Grosse K., Papernot N., Manoharan P., Backes M., McDaniel P. Adversarial Examples for Malware Detection. \emph{Computer Security --- ESORICS 2017}. LNCS Vol.~10493. Cham: Springer, 2017. P.~62--79.

\bibitem{ch01:grgof} Grégio A. R. A. et al. Toward a Taxonomy of Malware Behaviors. \emph{The Computer Journal}. 2015. Vol.~58, No.~10. P.~2758--2777. DOI: 10.1093/comjnl/bxv047.

\bibitem{ch13:gtfobins} GTFOBins Project. Curated List of Unix Binaries That Can Be Used to Bypass Local Security Restrictions. URL: \url{https://gtfobins.github.io/}.

\bibitem{ch08:guo} Guo F., Ferrie P., Chiueh T. A Study of the Packer Problem and Its Solutions. \emph{Recent Advances in Intrusion Detection (RAID 2008)}. LNCS Vol.~5230. Berlin: Springer, 2008. P.~98--115.

\bibitem{ch05:hoglund} Hoglund G., Butler J. Rootkits: Subverting the Windows Kernel. Boston: Addison-Wesley, 2005. 352~p.

\bibitem{ch05:barnea} Hosseini A. Ten Process Injection Techniques: A Technical Survey of Common and Trending Process Injection Techniques. Elastic Security Research, 2017.

\bibitem{lockheed-kc} Hutchins E. M., Cloppert M. J., Amin R. M. Intelligence-Driven Computer Network Defense Informed by Analysis of Adversary Campaigns and Intrusion Kill Chains. \emph{Leading Issues in Information Warfare \& Security Research}. 2011. Vol.~1. P.~80--106.

\bibitem{ch14:iec62443} International Electrotechnical Commission. IEC~62443: Industrial Communication Networks --- Network and System Security. Geneva: IEC, 2018--2021.

\bibitem{ch17:iso27001} International Organization for Standardization. ISO/IEC~27001: Information Security Management Systems --- Requirements. Geneva: ISO, 2022.

\bibitem{ch16:iso27037} International Organization for Standardization. ISO/IEC~27037: Guidelines for Identification, Collection, Acquisition and Preservation of Digital Evidence. Geneva: ISO, 2012.

\bibitem{iso-22301} ISO 22301:2019. Security and resilience --- Business continuity management systems --- Requirements. Geneva: ISO, 2019.

\bibitem{iso-31000} ISO 31000:2018. Risk management --- Guidelines. Geneva: ISO, 2018. URL: \url{https://www.iso.org/standard/65694.html}.

\bibitem{iso-27005} ISO/IEC 27005:2022. Information security, cybersecurity and privacy protection --- Guidance on managing information security risks. Geneva: ISO, 2022. URL: \url{https://www.iso.org/standard/80585.html}.

\bibitem{ch02:iso27032} ISO/IEC 27032:2023. Cybersecurity --- Guidelines for Internet security. Geneva: ISO, 2023. URL: \url{https://www.iso.org/standard/76070.html}.

\bibitem{iso-27035} ISO/IEC 27035-1:2023. Information technology --- Information security incident management --- Part~1: Principles and process. Geneva: ISO, 2023. URL: \url{https://www.iso.org/standard/78973.html}.

\bibitem{iso-27037} ISO/IEC 27037:2012. Guidelines for identification, collection, acquisition and preservation of digital evidence. Geneva: ISO, 2012. URL: \url{https://www.iso.org/standard/44381.html}.

\bibitem{ch06:jang} Jang J., Brumley D., Venkataraman S. BitShred: Feature Hashing Malware for Scalable Triage and Semantic Analysis. \emph{Proceedings of the 18th ACM Conference on Computer and Communications Security (CCS~'11)}. New York: ACM, 2011. P.~309--320.

\bibitem{ch08:junod} Junod P., Rinaldini J., Wehrli J., Michielin J. Obfuscator-LLVM --- Software Protection for the Masses. \emph{2015 IEEE/ACM 1st International Workshop on Software Protection (SPRO)}. Florence: IEEE, 2015. P.~3--9.

\bibitem{ch10:mitre-edr} Karantzas G., Patsakis C. An Empirical Assessment of Endpoint Detection and Response Systems against Advanced Persistent Threats Attack Vectors. \emph{Journal of Cybersecurity and Privacy}. 2021. Vol.~1, №~3. P.~387--421.

\bibitem{ch03:kay} Kay S. M. Fundamentals of Statistical Signal Processing. Vol.~II: Detection Theory. Upper Saddle River: Prentice Hall, 1998. 560~p.

\bibitem{ch07:kirat} Kirat D., Vigna G., Kruegel C. BareCloud: Bare-metal Analysis-based Evasive Malware Detection. \emph{Proceedings of the 23rd USENIX Security Symposium}. USENIX Association, 2014. P.~287--301.

\bibitem{ch05:klein} Klein T., Kotler I. Process Injection Techniques: Gotta Catch Them All. \emph{Black Hat USA}. Las Vegas, 2019.

\bibitem{ch03:kolter} Kolter J. Z., Maloof M. A. Learning to Detect and Classify Malicious Executables in the Wild. \emph{Journal of Machine Learning Research}. 2006. Vol.~7. P.~2721--2744.

\bibitem{ch04:kordy} Kordy B., Piètre-Cambacédès L., Schweitzer P. DAG-based Attack and Defense Modeling: Don't Miss the Forest for the Attack Trees. \emph{Computer Science Review}. 2014. Vol.~13--14. P.~1--38. DOI: 10.1016/j.cosrev.2014.07.001.

\bibitem{ch03:kornblum} Kornblum J. Identifying Almost Identical Files Using Context Triggered Piecewise Hashing. \emph{Digital Investigation}. 2006. Vol.~3 (Supplement). P.~91--97. DOI: 10.1016/j.diin.2006.06.015.

\bibitem{ch11:fileless} Kumar S., et al. An Emerging Threat: Fileless Malware --- A Survey and Research Challenges. \emph{Cybersecurity}. 2020. Vol.~3. Art.~1.

\bibitem{ch15:reproducible} Lamb C., Zacchiroli S. Reproducible Builds: Increasing the Integrity of Software Supply Chains. \emph{IEEE Software}. 2022. Vol.~39, №~2. P.~62--70.

\bibitem{ch14:langner} Langner R. Stuxnet: Dissecting a Cyberwarfare Weapon. \emph{IEEE Security \& Privacy}. 2011. Vol.~9, №~3. P.~49--51. DOI: 10.1109/MSP.2011.67.

\bibitem{ch13:hunting-maturity} Lee R. M., Bianco D. The Hunting Maturity Model. SANS Institute / Threat Hunting Project, 2016.

\bibitem{ch13:ocm} Lee R. M., Lee R. The Who, What, Where, When, Why and How of Effective Threat Hunting. SANS Institute Reading Room, 2016.

\bibitem{ch03:li} Li M., Vitányi P. An Introduction to Kolmogorov Complexity and Its Applications. 3rd ed. New York: Springer, 2008. 790~p.

\bibitem{ch07:ligh} Ligh M. H., Case A., Levy J., Walters A. The Art of Memory Forensics: Detecting Malware and Threats in Windows, Linux, and Mac Memory. Indianapolis: Wiley, 2014. 912~p.

\bibitem{ch07:lindorfer} Lindorfer M., Kolbitsch C., Comparetti P. M. Detecting Environment-Sensitive Malware. \emph{Recent Advances in Intrusion Detection (RAID 2011)}. LNCS Vol.~6961. Berlin: Springer, 2011. P.~338--357.

\bibitem{ch11:lolbas} LOLBAS Project. Living Off The Land Binaries, Scripts and Libraries. URL: \url{https://lolbas-project.github.io/}.

\bibitem{ch03:lyda} Lyda R., Hamrock J. Using Entropy Analysis to Find Encrypted and Packed Malware. \emph{IEEE Security \& Privacy}. 2007. Vol.~5, No.~2. P.~40--45. DOI: 10.1109/MSP.2007.48.

\bibitem{ch13:mandiant} Mandiant. APT1: Exposing One of China's Cyber Espionage Units. Mandiant Intelligence Center Report, 2013.

\bibitem{ch05:imphash} Mandiant. Tracking Malware with Import Hashing (imphash). Mandiant Threat Research, 2014.

\bibitem{ch05:matrosov} Matrosov A., Rodionov E., Bratus S. Rootkits and Bootkits: Reversing Modern Malware and Next Generation Threats. San Francisco: No Starch Press, 2019. 504~p.

\bibitem{ch12:mcmillan} McMillan R. Definition: Threat Intelligence. Gartner Research, 2013.

\bibitem{ch14:cloud-nist} Mell P., Grance T. The NIST Definition of Cloud Computing. NIST Special Publication 800-145. Gaithersburg: NIST, 2011.

\bibitem{ch10:amsi} Microsoft. Antimalware Scan Interface (AMSI) --- Technical Documentation. Microsoft Docs, 2024. URL: \url{https://learn.microsoft.com/windows/win32/amsi/antimalware-scan-interface-portal}.

\bibitem{ch10:etw} Microsoft. Event Tracing for Windows (ETW) --- Technical Documentation. Microsoft Docs, 2024. URL: \url{https://learn.microsoft.com/windows/win32/etw/about-event-tracing}.

\bibitem{ch05:pecoff} Microsoft. PE Format Specification (Portable Executable and Common Object File Format). Microsoft Docs, 2024. URL: \url{https://learn.microsoft.com/windows/win32/debug/pe-format}.

\bibitem{ch02:misp} MISP: Open Source Threat Intelligence Platform \& Standards for Threat Information Sharing. URL: \url{https://www.misp-project.org/}.

\bibitem{ch05:attck-persist} MITRE ATT\&CK. Tactic TA0003: Persistence. The MITRE Corporation. URL: \url{https://attack.mitre.org/tactics/TA0003/}.

\bibitem{ch05:attck-t1055} MITRE ATT\&CK. Technique T1055: Process Injection. The MITRE Corporation. URL: \url{https://attack.mitre.org/techniques/T1055/}.

\bibitem{mitre-attack} MITRE ATT\&CK\textsuperscript{\textregistered}: Adversarial Tactics, Techniques, and Common Knowledge. The MITRE Corporation. URL: \url{https://attack.mitre.org/}.

\bibitem{mitre-defend} MITRE D3FEND\textsuperscript{\texttrademark}: A Knowledge Graph of Cybersecurity Countermeasures. The MITRE Corporation. URL: \url{https://d3fend.mitre.org/}.

\bibitem{mitre-car} MITRE. Cyber Analytics Repository (CAR). The MITRE Corporation. URL: \url{https://car.mitre.org/}.

\bibitem{ch16:sans-ir} Murray J. Incident Handler's Handbook. SANS Institute Reading Room, 2015.

\bibitem{ch12:wannacry} Nakashima E., Timberg C. NSA Officials Worried about the Day Its Potent Hacking Tool Would Get Loose. Then It Did. \emph{The Washington Post}. 2017. (аналіз кампанії WannaCry).

\bibitem{nist-csf2} National Institute of Standards and Technology. The NIST Cybersecurity Framework (CSF) 2.0. NIST CSWP~29. Gaithersburg: NIST, 2024. DOI: 10.6028/NIST.CSWP.29.

\bibitem{ch10:defense-depth} National Security Agency. Defense in Depth: A Practical Strategy for Achieving Information Assurance in Today's Highly Networked Environments. NSA, 2010.

\bibitem{ch15:sbom} National Telecommunications and Information Administration (NTIA). The Minimum Elements for a Software Bill of Materials (SBOM). U.S. Department of Commerce, 2021.

\bibitem{ch15:sigstore} Newman J., Torres-Arias S., et al. Sigstore: Software Signing for Everybody. \emph{Proceedings of the 2022 ACM SIGSAC Conference on Computer and Communications Security (CCS~'22)}. New York: ACM, 2022. P.~2353--2367.

\bibitem{ch04:nist30} NIST. Guide for Conducting Risk Assessments. NIST SP 800-30, Rev.~1. Gaithersburg, MD: NIST, 2012. DOI: 10.6028/NIST.SP.800-30r1.

\bibitem{stix} OASIS. STIX Version 2.1: Structured Threat Information Expression. OASIS Standard, 2021. URL: \url{https://docs.oasis-open.org/cti/stix/v2.1/stix-v2.1.html}.

\bibitem{ch12:taxii} OASIS. TAXII Version 2.1: Trusted Automated Exchange of Intelligence Information. OASIS Standard, 2021.

\bibitem{ch15:ohm} Ohm M., Plate H., Sykosch A., Meier M. Backstabber's Knife Collection: A Review of Open Source Software Supply Chain Attacks. \emph{Detection of Intrusions and Malware, and Vulnerability Assessment (DIMVA 2020)}. LNCS Vol.~12223. Cham: Springer, 2020. P.~23--43.

\bibitem{ch06:tlsh} Oliver J., Cheng C., Chen Y. TLSH --- A Locality Sensitive Hash. \emph{2013 Fourth Cybercrime and Trustworthy Computing Workshop}. Sydney: IEEE, 2013. P.~7--13.

\bibitem{ch11:suricata} Open Information Security Foundation (OISF). Suricata: Open Source Intrusion Detection and Prevention Engine --- User Guide. URL: \url{https://docs.suricata.io/}.

\bibitem{ch15:slsa} Open Source Security Foundation (OpenSSF). SLSA: Supply-chain Levels for Software Artifacts. URL: \url{https://slsa.dev/}.

\bibitem{ch14:iot-owasp} OWASP Foundation. OWASP Internet of Things (IoT) Top 10. OWASP Project, 2018.

\bibitem{sigma} Patzke T. Sigma: Generic Signature Format for SIEM Systems. SigmaHQ Project Documentation. URL: \url{https://github.com/SigmaHQ/sigma}.

\bibitem{ch04:pols} Pols P. The Unified Kill Chain: Designing a Unified Kill Chain for analyzing, comparing and defending against cyber attacks: Master thesis. Cyber Security Academy, 2017.

\bibitem{ch09:raff} Raff E., Barker J., Sylvester J., Brandon R., Catanzaro B., Nicholas C. Malware Detection by Eating a Whole EXE. \emph{AAAI Workshop on Artificial Intelligence for Cyber Security}. 2018.

\bibitem{gdpr} Regulation (EU) 2016/679 (General Data Protection Regulation). \emph{Official Journal of the EU}. L 119. P.~1--88. URL: \url{https://eur-lex.europa.eu/eli/reg/2016/679/oj}.

\bibitem{ch02:csa} Regulation (EU) 2019/881 on ENISA and on information and communications technology cybersecurity certification (Cybersecurity Act). \emph{Official Journal of the EU}. L 151. P.~15--69. URL: \url{https://eur-lex.europa.eu/eli/reg/2019/881/oj}.

\bibitem{ch02:eccc} Regulation (EU) 2021/887 establishing the European Cybersecurity Industrial, Technology and Research Competence Centre and the Network of National Coordination Centres. \emph{Official Journal of the EU}. L 202. 2021.

\bibitem{cra} Regulation (EU) 2024/2847 on horizontal cybersecurity requirements for products with digital elements (Cyber Resilience Act). \emph{Official Journal of the EU}. 2024. URL: \url{https://eur-lex.europa.eu/eli/reg/2024/2847/oj}.

\bibitem{ch03:rice} Rice H. G. Classes of Recursively Enumerable Sets and Their Decision Problems. \emph{Transactions of the American Mathematical Society}. 1953. Vol.~74, No.~2. P.~358--366. DOI: 10.1090/S0002-9947-1953-0053041-6.

\bibitem{ch12:rid} Rid T., Buchanan B. Attributing Cyber Attacks. \emph{Journal of Strategic Studies}. 2015. Vol.~38, №~1--2. P.~4--37. DOI: 10.1080/01402390.2014.977382.

\bibitem{ch11:snort} Roesch M. Snort --- Lightweight Intrusion Detection for Networks. \emph{Proceedings of the 13th USENIX Conference on System Administration (LISA '99)}. USENIX Association, 1999. P.~229--238.

\bibitem{ch17:resilience} Ross R., Pillitteri V., Graubart R., Bodeau D., McQuaid R. Developing Cyber-Resilient Systems: A Systems Security Engineering Approach. NIST Special Publication 800-160, Vol.~2, Rev.~1. Gaithersburg: NIST, 2021.

\bibitem{ch08:royal} Royal P., Halpin M., Dagon D., Edmonds R., Lee W. PolyUnpack: Automating the Hidden-Code Extraction of Unpack-Executing Malware. \emph{22nd Annual Computer Security Applications Conference (ACSAC)}. Miami Beach: IEEE, 2006. P.~289--300.

\bibitem{ch10:sysmon} Russinovich M., Garnier T. System Monitor (Sysmon) --- Windows Sysinternals Documentation. Microsoft, 2024. URL: \url{https://learn.microsoft.com/sysinternals/downloads/sysmon}.

\bibitem{ch05:russinovich} Russinovich M., Solomon D., Ionescu A., Allievi A. Windows Internals. 7th ed. Parts~1--2. Redmond: Microsoft Press, 2017--2021.

\bibitem{ch10:least-priv} Saltzer J. H., Schroeder M. D. The Protection of Information in Computer Systems. \emph{Proceedings of the IEEE}. 1975. Vol.~63, №~9. P.~1278--1308.

\bibitem{ch09:saxe} Saxe J., Berlin K. Deep Neural Network Based Malware Detection Using Two Dimensional Binary Program Features. \emph{2015 10th International Conference on Malicious and Unwanted Software (MALWARE)}. Fajardo: IEEE, 2015. P.~11--20.

\bibitem{ch04:schneier} Schneier B. Attack Trees. \emph{Dr. Dobb's Journal}. 1999. Vol.~24, No.~12. P.~21--29.

\bibitem{ch08:oreans} Schrittwieser E., Katzenbeisser S., Kinder J., Merzdovnik G., Weippl E. Protecting Software through Obfuscation: Can It Keep Pace with Progress in Code Analysis? \emph{ACM Computing Surveys}. 2016. Vol.~49, №~1. Art.~4. P.~1--37.

\bibitem{ch09:schultz} Schultz M. G., Eskin E., Zadok E., Stolfo S. J. Data Mining Methods for Detection of New Malicious Executables. \emph{2001 IEEE Symposium on Security and Privacy}. Oakland: IEEE, 2001. P.~38--49.

\bibitem{ch02:budapest-ap} Second Additional Protocol to the Convention on Cybercrime on enhanced co-operation and disclosure of electronic evidence. CETS No.~224. Strasbourg: Council of Europe, 2022.

\bibitem{ch10:appcontrol} Sedgewick A., Souppaya M., Scarfone K. Guide to Application Whitelisting. NIST Special Publication 800-167. Gaithersburg: NIST, 2015. DOI: 10.6028/NIST.SP.800-167.

\bibitem{ch03:selcuk} Selçuk A. A., Orhan F., Batur B. Undecidable Problems in Malware Analysis. \emph{2017 12th International Conference for Internet Technology and Secured Transactions (ICITST)}. Cambridge: IEEE, 2017. P.~494--497.

\bibitem{ch03:shannon} Shannon C. E. A Mathematical Theory of Communication. \emph{Bell System Technical Journal}. 1948. Vol.~27, No.~3. P.~379--423.

\bibitem{sikorski} Sikorski M., Honig A. Practical Malware Analysis: The Hands-On Guide to Dissecting Malicious Software. San Francisco: No Starch Press, 2012. 800~p.

\bibitem{industroyer2} Slowik J., Hultquist J. et al. Industroyer2: Nozomi Networks Labs Analyses the IEC~104 Payload: technical report. 2022. URL: \url{https://cert.gov.ua/article/39518}.

\bibitem{ch13:lateral} Smiliotopoulos C., Kambourakis G., Barbatsalou K. On the Detection of Lateral Movement through Supervised Machine Learning and an Open-Source Tool to Create Turnkey Datasets. \emph{International Journal of Information Security}. 2023. Vol.~22. P.~1893--1919.

\bibitem{ch09:sommer} Sommer R., Paxson V. Outside the Closed World: On Using Machine Learning for Network Intrusion Detection. \emph{2010 IEEE Symposium on Security and Privacy}. Oakland: IEEE, 2010. P.~305--316.

\bibitem{ch14:container} Souppaya M., Morello J., Scarfone K. Application Container Security Guide. NIST Special Publication 800-190. Gaithersburg: NIST, 2017.

\bibitem{ch16:recovery} Souppaya M., Scarfone K. Guide for Cybersecurity Event Recovery. NIST Special Publication 800-184. Gaithersburg: NIST, 2016.

\bibitem{nist-83} Souppaya M., Scarfone K. Guide to Malware Incident Prevention and Handling for Desktops and Laptops. NIST Special Publication 800-83, Revision~1. Gaithersburg: NIST, 2013. DOI: 10.6028/NIST.SP.800-83r1.

\bibitem{ch01:spafford} Spafford E. H. The Internet Worm Program: An Analysis. \emph{ACM SIGCOMM Computer Communication Review}. 1989. Vol.~19, No.~1. P.~17--57. DOI: 10.1145/66093.66095.

\bibitem{ch03:spinellis} Spinellis D. Reliable Identification of Bounded-length Viruses is NP-complete. \emph{IEEE Transactions on Information Theory}. 2003. Vol.~49, No.~1. P.~280--284. DOI: 10.1109/TIT.2002.806137.

\bibitem{ch13:sqrrl} Sqrrl Team. A Framework for Cyber Threat Hunting. Sqrrl Data, 2016.

\bibitem{ch04:staniford} Staniford S., Paxson V., Weaver N. How to 0wn the Internet in Your Spare Time. \emph{Proceedings of the 11th USENIX Security Symposium}. USENIX Association, 2002. P.~149--167.

\bibitem{ch11:detection-maturity} Stillions R. The DML Model (Detection Maturity Level). ryan.stillions blog, 2014.

\bibitem{ch13:stoll} Stoll C. The Cuckoo's Egg: Tracking a Spy Through the Maze of Computer Espionage. New York: Doubleday, 1989. 326~p.

\bibitem{ch14:stouffer} Stouffer K., Pillitteri V., Lightman S., Abrams M., Hahn A. Guide to Industrial Control Systems (ICS) Security. NIST Special Publication 800-82, Revision~2. Gaithersburg: NIST, 2015. DOI: 10.6028/NIST.SP.800-82r2.

\bibitem{ch04:strom} Strom B. E. et al. MITRE ATT\&CK: Design and Philosophy. MITRE Technical Report MTR190121. McLean: The MITRE Corporation, 2020.

\bibitem{szor} Szor P. The Art of Computer Virus Research and Defense. Boston: Addison-Wesley, 2005. 744~p.

\bibitem{ch02:taxii} TAXII Version 2.1: Trusted Automated Exchange of Intelligence Information. OASIS Standard, 2021. URL: \url{https://docs.oasis-open.org/cti/taxii/v2.1/}.

\bibitem{ch07:volatility} The Volatility Foundation. The Volatility Framework: Volatile Memory Artifact Extraction Utility Framework. URL: \url{https://www.volatilityfoundation.org/}.

\bibitem{ch05:elf} Tool Interface Standard (TIS). Executable and Linkable Format (ELF) Specification. Version~1.2. TIS Committee, 1995.

\bibitem{ch15:intoto} Torres-Arias S., Afzali H., Kuppusamy T. K., Curtmola R., Cappos J. in-toto: Providing Farm-to-Table Guarantees for Bits and Bytes. \emph{Proceedings of the 28th USENIX Security Symposium}. USENIX Association, 2019. P.~1393--1410.

\bibitem{ch17:c2m2} U.S. Department of Energy. Cybersecurity Capability Maturity Model (C2M2), Version~2.1. Washington: DOE, 2022.

\bibitem{ch09:ucci} Ucci D., Aniello L., Baldoni R. Survey of Machine Learning Techniques for Malware Analysis. \emph{Computers \& Security}. 2019. Vol.~81. P.~123--147. DOI: 10.1016/j.cose.2018.11.001.

\bibitem{ch08:ugarte} Ugarte-Pedrero X., Balzarotti D., Santos I., Bringas P. G. SoK: Deep Packer Inspection --- A Longitudinal Study of the Complexity of Run-Time Packers. \emph{2015 IEEE Symposium on Security and Privacy}. San Jose: IEEE, 2015. P.~659--673. DOI: 10.1109/SP.2015.46.

\bibitem{ch01:dbir} Verizon. Data Breach Investigations Report (DBIR). New York: Verizon Business, 2024. URL: \url{https://www.verizon.com/business/resources/reports/dbir/}.

\bibitem{ch17:verizon} Verizon. Data Breach Investigations Report (DBIR). Verizon Business, 2024. (щорічний аналітичний звіт про чинник людини в інцидентах).

\bibitem{ch13:virvilis} Virvilis N., Gritzalis D. The Big Four --- What We Did Wrong in Advanced Persistent Threat Detection? \emph{2013 International Conference on Availability, Reliability and Security}. Regensburg: IEEE, 2013. P.~248--254.

\bibitem{ch01:vonneumann} von Neumann J. Theory of Self-Reproducing Automata / ed. A.~W.~Burks. Urbana: University of Illinois Press, 1966.

\bibitem{ch12:misp} Wagner C., Dulaunoy A., Wagener G., Iklody A. MISP --- The Design and Implementation of a Collaborative Threat Intelligence Sharing Platform. \emph{Proceedings of the 2016 ACM on Workshop on Information Sharing and Collaborative Security (WISCS~'16)}. New York: ACM, 2016. P.~49--56.

\bibitem{ch07:willems} Willems C., Holz T., Freiling F. Toward Automated Dynamic Malware Analysis Using CWSandbox. \emph{IEEE Security \& Privacy}. 2007. Vol.~5, №~2. P.~32--39.

\bibitem{ch08:yadegari} Yadegari B., Johannesmeyer B., Whitely B., Debray S. A Generic Approach to Automatic Deobfuscation of Executable Code. \emph{2015 IEEE Symposium on Security and Privacy}. San Jose: IEEE, 2015. P.~674--691.

\bibitem{ch08:you} You I., Yim K. Malware Obfuscation Techniques: A Brief Survey. \emph{2010 International Conference on Broadband, Wireless Computing, Communication and Applications}. Fukuoka: IEEE, 2010. P.~297--300.

\bibitem{ch01:zetter} Zetter K. Countdown to Zero Day: Stuxnet and the Launch of the World's First Digital Weapon. New York: Crown, 2014. 448~p.

\bibitem{ch11:alerting} Zimmerman C. Ten Strategies of a World-Class Cybersecurity Operations Center. McLean: MITRE Corporation, 2014. 353~p.

\bibitem{ch04:zou} Zou C. C., Towsley D., Gong W. Email Worm Modeling and Defense. \emph{Proceedings of the 13th International Conference on Computer Communications and Networks (ICCCN)}. IEEE, 2004. P.~409--414.

\bibitem{ch06:yara-doc} Álvarez V. M. YARA Documentation: Writing YARA rules. Release~4.x. VirusTotal. URL: \url{https://yara.readthedocs.io/}.

\bibitem{yara} Álvarez V. M. YARA: The pattern matching swiss knife for malware researchers. URL: \url{https://virustotal.github.io/yara/}.
